\documentclass[11pt,a4paper]{article}
\usepackage[T1]{fontenc}
\usepackage{amsmath,amssymb,amsthm,mathtools}
\usepackage{newpxtext,newpxmath}
\usepackage[margin=27mm,headheight=14pt]{geometry}
\usepackage{microtype,booktabs,longtable,enumitem,fancyhdr}
\usepackage[hidelinks]{hyperref}
\setlength{\emergencystretch}{2em}
\setlength{\parskip}{.25em}
\newtheorem{theorem}{Theorem}[section]
\newtheorem{lemma}[theorem]{Lemma}
\newtheorem{corollary}[theorem]{Corollary}
\newtheorem{proposition}[theorem]{Proposition}
\theoremstyle{definition}\newtheorem{definition}[theorem]{Definition}
\theoremstyle{remark}\newtheorem*{remark}{Remark}
\newcommand{\N}{\mathbb N}\newcommand{\Q}{\mathbb Q}\newcommand{\Z}{\mathbb Z}\newcommand{\C}{\mathbb C}
\newcommand{\ph}{\varphi}\newcommand{\Span}{\operatorname{span}}\newcommand{\ord}{\operatorname{ord}}
\newcommand{\sig}{\sigma}\newcommand{\tauu}{\tau}\newcommand{\rad}{\operatorname{rad}}
\pagestyle{fancy}\fancyhf{}\fancyhead[L]{\small Affine sections and dilation closure}\fancyhead[R]{\small\thepage}\renewcommand{\headrulewidth}{0pt}
\title{Affine Sections, Local Factors, and Dilation Closure\\[4pt]\large Round-five research memorandum for Erd\H{o}s \#249}
\author{Research candidates for independent review}
\date{7 September 2026}
\begin{document}\maketitle
\begin{abstract}
Restriction to arithmetic progressions separates variable coefficients from affine arithmetic functions. This gives a shorter proof of the existing affine-totient signature classification and removes its bespoke isolation hypothesis from the ordinary argument. The same principle yields proposed exact all-base kernel ranks for the sum-of-divisors and divisor-counting functions. In a base with two prime factors, the ranks differ by one from depth three onwards: their local translation spaces are exponential and polynomial, respectively. A second theorem characterises finite rational dilation closure for bounded integer-weighted Mersenne series. Closure holds exactly for eventually periodic weights; every other weight produces a natural boundary. The M\"obius specialisation therefore excludes finite rational linear dilation closure, although its first-power value is rational. Complete ordinary proofs and exact finite checks accompany the statements. None is asserted to have been compiled in Lean, independently reviewed, or established as a priority claim. No irrationality proof for the unreduced totient series is supplied.
\end{abstract}

\section{Separate arithmetic progressions before separating coefficients}

The current paper proves rational linear independence of
$\ph(a_i n+b_i)$ for positive slopes and pairwise nonproportional affine forms.
Martin's stronger ordering theorem is an antecedent for this independence, and
his Corollary~4 supplies the corresponding statement for $\sigma$~\cite{martin}.
The current Lean proof uses a diagonal evaluation matrix instead.
The following elementary transport should be used before constructing another
prime-isolation argument.

\begin{definition}
A finite family $w_i:\N_{>0}\to\Q$ can be \emph{frozen on progressions} if, for
every $n_0\ge1$, there is $Q\ge1$ such that
\[
 w_i(n_0+Qt)=w_i(n_0)\qquad(i,\ t\ge0).
\]
The modulus may depend on $n_0$. A common period is sufficient but unnecessary.
\end{definition}

\begin{lemma}[Progression freezing]\label{lem:freeze}
Let $L_i(n)=a_i n+b_i$, with $a_i\ge1$, $b_i\ge0$, and
$a_i b_j\ne a_j b_i$ for $i\ne j$. Suppose an arithmetic function $f$ has
independent values on every such finite family of affine forms, including
after restriction to an arithmetic progression. If $w_i$ can be frozen on
progressions and
\[
 \sum_i w_i(n)f(L_i(n))=0\qquad(n\ge1),
\]
then every $w_i$ vanishes.
\end{lemma}
\begin{proof}
Fix $n_0$ and a freezing modulus $Q$. Substitute $n=n_0+Qt$.
The coefficients become the constants $w_i(n_0)$, and the new forms have slopes
$a_iQ$ and intercepts $a_i n_0+b_i>0$. Their cross determinants are
\[
 (a_iQ)(a_jn_0+b_j)-(a_jQ)(a_in_0+b_i)
 =Q(a_i b_j-a_j b_i)\ne0.
\]
Affine independence gives $w_i(n_0)=0$. Since $n_0$ was arbitrary, this proves
the claim. The same argument applies to an eventual relation at every $n_0$
beyond its threshold. For periodic coefficients, shifting each residue class
past that threshold proves vanishing everywhere.
\end{proof}

\begin{corollary}[Periodic coefficient independence]\label{cor:periodic}
Pairwise nonproportional affine totient sequences are independent over the
ring of rational-valued periodic sequences. The same assertion holds for
$\sigma$ and $\tau$.
\end{corollary}
\begin{proof}
Take a common period and apply Lemma~\ref{lem:freeze}. Independence for $\tau$
is proved below; the other two cases are also consequences of Martin's
results. This is a ring-independence assertion about sequences, not about
values of their generating series.
\end{proof}

Functions of finitely many exact valuations are an important nonperiodic
example. Given finitely many positive affine values $L_j(n_0)$ and primes $p$,
choose $Q$ divisible by
\[
 p^{1+\max_j v_p(L_j(n_0))}
\]
for every relevant $p$. Then every $v_p(L_j(n_0+Qt))$ is unchanged for $t\ge0$.
No uniform modulus over all $n_0$ is being claimed.

\subsection{A constructive inverse for the existing signature theorem}

For $L(n)=An+B$ primitive, $A\ge1$, $B\ge0$, and $c\ge1$, put
\[
 P(c,A)=\{p:p\mid c,\ p\nmid A\}.
\]
The local totient identity gives, at positive arguments,
\begin{equation}\label{eq:content}
 \frac{\ph(cL(n))}{\ph(c)}
 =\ph(L(n))\prod_{\substack{p\in P(c,A)\\p\mid L(n)}}\frac p{p-1}.
\end{equation}
The multiplier on the right is periodic. A relation among mixed primitive
forms therefore separates into one relation per primitive form by
Corollary~\ref{cor:periodic}. This step requires no new isolation hypothesis.

Fix one such form and let $P$ be the union of its variable content primes.
Every pattern $T\subseteq P$ occurs as $\{p\in P:p\mid L(n)\}$: $A$ is a unit
modulo every $p\in P$, so the Chinese remainder theorem prescribes zero or a
nonzero residue independently. For sets $S,T\subseteq P$, define
\[
 K(T,S)=\prod_{p\in T\cap S}\rho_p,\qquad \rho_p=\frac p{p-1}.
\]
Indexing by the full Boolean cube makes the inverse explicit:
\begin{equation}\label{eq:tensor}
 K=\bigotimes_{p\in P}
 \begin{pmatrix}1&1\\1&\rho_p\end{pmatrix},\qquad
 K^{-1}=\bigotimes_{p\in P}
 \begin{pmatrix}p&-(p-1)\\-(p-1)&p-1\end{pmatrix}.
\end{equation}
Both assertions follow by multiplying the two $2\times2$ factors.
In particular the inverse is integral. For $h=|P|\ge1$,
\[
 \det K=\prod_{p\in P}(p-1)^{-2^{h-1}}.
\]
Distinct selected columns remain independent by extending their coefficients
by zero to the full cube. The positive-definiteness argument in the previous
round is thus replaced by coordinate recovery. The general meet-matrix and
incidence-factorisation framework is classical~\cite{mattila}; the useful
specialisation here is the integral inverse and the order of the arithmetic
reduction.

\begin{theorem}[Signature completeness, with a shorter proof]\label{thm:signature}
For a finite family $a_i\ge1$, $b_i\ge0$, set
\[
 c_i=\gcd(a_i,b_i),\quad A_i=a_i/c_i,\quad B_i=b_i/c_i,
 \quad \Sigma_i=(A_i,B_i,P(c_i,A_i)).
\]
Then
\[
 \sum_i\lambda_i\ph(a_i n+b_i)=0\quad(n\ge0)
 \quad\Longleftrightarrow\quad
 \sum_{i:\Sigma_i=\Sigma}\lambda_i\ph(c_i)=0\quad\text{for every }\Sigma.
\]
The conclusion is unchanged when the relation is required only eventually.
The rational dimension of the family is the number of distinct signatures.
\end{theorem}
\begin{proof}
Group by primitive form and use \eqref{eq:content}. Progression freezing
separates these groups. For one group, aggregate coefficients with the same
set $S=P(c_i,A_i)$. Evaluation on all divisibility patterns gives $K\lambda=0$;
\eqref{eq:tensor} gives $\lambda=0$. The converse is the content identity.
Every pattern occurs arbitrarily late, so the argument also treats eventual
relations. The only possible zero argument is $L(n)=n$ at $n=0$; all its
content multiples have totient zero there, so that endpoint causes no change.
\end{proof}

This theorem was already an ordinary-proof result in round four. The new
contribution in this section is a shorter proof and a smaller formalisation
interface. The attached \texttt{PeriodicTotientIndependence.lean} and
\texttt{BooleanEulerInverse.lean} are uncompiled candidates for those two
interfaces, not a completed Lean proof of Theorem~\ref{thm:signature}.

\subsection{Two inexpensive consequences for the current flagship}

\begin{proposition}[Integral normal form and its exact scope]\label{prop:integral}
For $k\ge2,e\ge1$, the current canonical totient family is a $\Z$-basis of the
$\Z$-module generated by the sections through level $e$. The integral relation
module has a basis consisting of one elementary scalar reduction per omitted
section. Its rank is $\sum_{j=1}^{e-1}k^j$.
\end{proposition}
\begin{proof}
For a nonzero omitted residue $r=k^tu$, $k\nmid u$, the reduction scalar is
\[
 C_k(t,u)=\frac{k^t}{\prod_{p\mid k,\ p\nmid u}p}
             \prod_{p\mid k,\ p\nmid u}(p-1)\in\Z.
\]
Zero-residue reductions also have integral scalars. Rational independence
therefore implies integral independence and integral spanning. In the free
module on all channels, each relation ``omitted channel minus its scalar
multiple of the retained channel'' has a unique omitted-coordinate pivot
$1$. Subtracting these relations eliminates all omitted coordinates, after
which independence eliminates the rest. Counting pivots gives the formula.
\end{proof}

The module in this proposition is not saturated in the full group of
integer-valued sequences. For $e\ge2$ at base $2$,
$\ph(4n+3)/2$ is integer-valued, but its unique coefficient on the retained
channel $\ph(4n+3)$ is $1/2$. It is outside the generated integral module.
This falsifies an ambient-saturation upgrade.

The rank and basis also survive deletion of any finite prefix. The affine
matrix rows can be chosen arbitrarily late. For the two zero channels, use
large $n\equiv1\pmod k$ and large $n=k^t$, giving the distinct ratios
$\ph(k)$ and $k$. This supplies an eventual-relation normal form without a
separate compactness argument.

\section{Exact all-base ranks for divisor functions}\label{sec:divisor}

Write $\sigma(n)=\sum_{d\mid n}d$ and $\tau(n)=\sum_{d\mid n}1$, and set both
functions to zero at $n=0$. For either $f$, define
\[
 V^f_{k,e}=\Span_\Q\{n\mapsto f(k^jn+r):0\le j\le e,\ 0\le r<k^j\}.
\]
The supplied corpus contains measured ranks $41$ for both functions at
$(k,e)=(6,2)$. These measurements are not being re-announced as a discovery.
The proposed result is a formula and a canonical rational basis at every
base and depth.

\subsection{Arithmetic independence of the primitive classes}

\begin{lemma}\label{lem:divisor-affine}
For $f=\sigma$ or $f=\tau$, a finite family $n\mapsto f(a_i n+b_i)$ with
positive slopes and pairwise nonproportional forms is rationally independent.
The evaluation points witnessing independence may be chosen arbitrarily late.
\end{lemma}
\begin{proof}
A finite shift makes all intercepts positive. Write $L_i=g_iA_i$, where $A_i$
is primitive. Choose an odd prime $\ell$ avoiding all slopes, contents,
nonzero cross determinants, and all nonzero integers $f(g_i)$.

For $\sigma$, fix row $i$. For each $j\ne i$, choose a fresh prime
$q_{ij}\equiv-1\pmod\ell$, avoiding the same finite set, and prescribe
\[
 L_j(n)\equiv q_{ij}\pmod{q_{ij}^2}.
\]
This forces $v_{q_{ij}}(L_j(n))=1$, so $\ell\mid\sigma(L_j(n))$.
Also prescribe $A_i(n)\equiv1\pmod\ell$. The CRT combines these conditions.
On the resulting progression, the primitive form $A_i(n)$ lies in a reduced
prime progression. Indeed it is a unit modulo each $q_{ij}$, since otherwise
the cross determinant for $L_i,L_j$ would vanish modulo $q_{ij}$; it is a unit
modulo $\ell$ by construction; and it is automatically coprime to its own
slope. Dirichlet supplies a prime $p=A_i(n)$, arbitrarily large and outside
$g_i$. The diagonal value is $\sigma(g_i)(p+1)\not\equiv0\pmod\ell$.

For $\tau$, choose fresh primes $q_{ij}$ avoiding the finite exceptional set,
and prescribe
\[
 L_j(n)\equiv q_{ij}^{\ell-1}\pmod{q_{ij}^{\ell}}.
\]
The exact valuation is $\ell-1$, so $\ell\mid\tau(L_j(n))$. The same
reducedness argument lets $A_i(n)=p$ be prime and outside $g_i$; now the
diagonal is $2\tau(g_i)\not\equiv0\pmod\ell$. No distribution theorem for
simultaneous prime values is used: only one selected primitive form is made
prime per row.

One such evaluation for each $i$ gives a matrix diagonal and nonsingular
modulo $\ell$. Its rational determinant is nonzero. For $\sigma$, Martin's
Corollary~4 is a stronger antecedent; the direct proof is included to display
the exact hypotheses needed for the transport argument.
\end{proof}

\subsection{The local dimensions}

Let $k=\prod_{p\mid k}p^{a_p}$. Each nonzero proportionality class of channels
has a unique first representative
\[
 (t,u),\qquad 1\le t\le e,\quad 1\le u<k^t,\quad k\nmid u.
\]
Its channels are $(t+s,k^su)$ for $0\le s\le e-t$. Put
\[
 g=\gcd(k^t,u),\qquad L(n)=\frac{k^t}{g}n+\frac ug,
\]
\begin{equation}\label{eq:Ptu}
 P_{t,u}=\{p\mid k:v_p(u)\ge t a_p\},\qquad h_{t,u}=|P_{t,u}|.
\end{equation}
These are exactly the primes of $k$ which do not divide the slope of $L$.
They are the primes whose valuations in $L(n)$ can vary.

\begin{lemma}[Local translation dimension]\label{lem:localdimension}
A class with $m$ consecutive content levels and $h$ variable primes has rank
\[
 d_\sigma(m,h)=\min(m,2^h),\qquad
 d_\tau(m,h)=\min(m,h+1).
\]
Its first $d_f(m,h)$ channels form a rational basis.
\end{lemma}
\begin{proof}
At primes of $k$ dividing the primitive slope, $L(n)$ is always a unit. Their
local factors are nonzero scalars depending only on $s$; divide each channel
by those scalars. At the variable primes, write $v_p=v_p(L(n))$ and
$c_p=v_p(g)$. After division by the common factor $f(L(n))$ and multiplication
by the common nonzero local denominators, the $\sigma$ functions become
\[
 H_s(v)=\prod_{p\in P}
 \frac{p^{v_p+c_p+1}p^{a_p s}-1}{p-1}
 =\sum_{T\subseteq P}C_T
       \left(\prod_{p\in T}p^{v_p}\right)\lambda_T^s,
 \qquad \lambda_T=\prod_{p\in T}p^{a_p},
\]
where every $C_T\ne0$. Unique prime factorisation makes the $2^h$ numbers
$\lambda_T$ distinct. The functions $\prod_{p\in T}p^{v_p}$ are independent:
restrict each $v_p$ to $0,1$ and use the tensor product of the invertible
matrices $\left(\begin{smallmatrix}1&1\\1&p\end{smallmatrix}\right)$.
All these valuation patterns are realised by CRT, with exact valuations
prescribed modulo $p^{v_p+1}$. A Vandermonde matrix in $s=0,\ldots,m-1$
therefore has rank $\min(m,2^h)$, and its first $\min(m,2^h)$ columns are
independent.

For $\tau$ the corresponding functions are
\[
 H_s(v)=\prod_{p\in P}(v_p+c_p+1+a_p s)
       =\sum_{j=0}^h E_j(v)s^j.
\]
The polynomial $E_j$ has total degree $h-j$ and a nonzero highest homogeneous
part. Thus $E_0,\ldots,E_h$ are independent polynomials. A polynomial
vanishing on all $v\in\N^h$ is zero, by induction on $h$, so they are also
independent on the realised valuation vectors. The Vandermonde matrix in
$s$ has rank $\min(m,h+1)$, with the asserted first-column basis.

For the zero class $L(n)=n$, the same proof applies on $n\ge1$ with $g=1$,
$c_p=0$, and $P$ the full prime support of $k$. At $n=0$ all channels vanish.
\end{proof}

\begin{theorem}[Proposed all-base divisor-kernel formula]\label{thm:divisor-rank}
For $k\ge2$, $e\ge0$, and $f\in\{\sigma,\tau\}$,
\begin{equation}\label{eq:rankformula}
 \dim_\Q V^f_{k,e}
 =d_f(e+1,\omega(k))+
 \sum_{t=1}^e\ \sum_{\substack{1\le u<k^t\\ k\nmid u}}
       d_f(e-t+1,h_{t,u}).
\end{equation}
A rational basis consists of the first $d_f(e+1,\omega(k))$ zero channels and,
in every nonzero class $(t,u)$, its first $d_f(e-t+1,h_{t,u})$ channels.
\end{theorem}
\begin{proof}
Every channel belongs to one of the specified proportionality classes:
remove the largest common power of $k$ from its nonzero residue. Two distinct
first representatives have different primitive forms, because equality of
the fractions $u/k^t$ would force the representative at the larger level to
have residue divisible by $k$.

Consider a relation among all classes. In a class with primitive form $L$,
factor out $f(L(n))$. The coefficient left over is a rational function of
finitely many exact valuations $v_p(L(n))$. At each $n_0\ge1$ all such
coefficients, for all classes simultaneously, can be frozen on one
progression. Lemmas~\ref{lem:freeze} and \ref{lem:divisor-affine} then force
each primitive-class coefficient to vanish separately. Hence dimensions add
across the classes. Lemma~\ref{lem:localdimension} supplies each summand and
its first-level basis. All within-class identities extend to $n=0$; the zero
class vanishes there and positive-intercept classes have positive values.
\end{proof}

After division by the scalar factors at primes dividing the primitive slope,
the class sequence in its content-level index $s$ has annihilator
\[
 \prod_{T\subseteq P}(X-\lambda_T)\quad\text{for }\sigma,
 \qquad (X-1)^{h+1}\quad\text{for }\tau.
\]
These give exact coordinate recurrences for every omitted channel. Unlike the
totient case, the displayed first-level basis need not be integral:
\[
 \sigma(64n+16)=\frac{31}{3}\sigma(8n+2).
\]
An assertion of integral spanning by that chosen basis would be false.

\begin{corollary}[Prime-power and two-prime-factor bases]\label{cor:twoprimes}
At a prime-power base $k$, both ranks are $k^e+1$ for $e\ge1$.
If $k=xy$ with $x=p^a$, $y=q^b$ for distinct primes $p,q$, then, for $e\ge1$,
\begin{align*}
 \dim V^\sigma_{k,e}
 &=k^e+x^{e-1}+y^{e-1}+\min(e+1,4)-3,\\
 \dim V^\tau_{k,e}
 &=k^e+x^{e-1}+y^{e-1}+\min(e+1,3)-3.
\end{align*}
\end{corollary}
\begin{proof}
For a prime-power base, $k\nmid u$ excludes the sole prime from $P_{t,u}$.
Every nonzero class has dimension one, and the zero class has dimension two.

For two prime factors, a nonzero class has at most one variable prime, since
$P_{t,u}$ containing both would imply $k^t\mid u$. At first level $t$, the
numbers of classes having the first or the second variable prime are
$(y-1)y^{t-1}$ and $(x-1)x^{t-1}$, respectively. Indeed, to have the first
prime variable one needs $x^t\mid u$, and $k\nmid u$ then excludes divisibility
by $y$; count the multiples of $x^t$ below $x^ty^t$ accordingly. Each such
class contributes one extra dimension exactly when $t\le e-1$.
The total nonzero first representatives number $k^e-1$. Summing the two
geometric progressions gives the formulas, with the appropriate zero-class
capacity.
\end{proof}

In particular, the successive ranks in base $6$ are
\[
\begin{array}{c|rrr}
 e&1&2&3\\\hline
 \sigma&7&41&230\\
 \tau&7&41&229\\
 \ph&7&37&217
\end{array}
\]
The distinction is local. With two variable primes, the sum-of-divisors
factors have four distinct exponential modes, while the divisor-counting
factors give a quadratic polynomial with three modes.

\begin{remark}[Positive divisor powers]
The same $\sigma$ rank formula holds for $\sigma_a(n)=\sum_{d\mid n}d^a$ at
every positive integer $a$. In the affine independence proof, choose
$\ell\equiv1\pmod{2a}$ outside the finite exceptional set and choose
$q_{ij}$ in a residue of order $2a$ modulo $\ell$. Exact valuation one then
gives $\sigma_a(q_{ij})\equiv0\pmod\ell$. A selected prime
$p\equiv1\pmod\ell$ gives $\sigma_a(p)\equiv2$. In the local dimension
proof replace $\lambda_T$ by $\lambda_T^a$, which remain distinct. At the
formal limit $a\to0$, each local divisor sum becomes its exponent plus one;
coalescence of the exponential modes explains the polynomial $\tau$ case.
This last limit is explanatory; the integer $a=0$ theorem was proved directly.
\end{remark}

\subsection{A rational generating function for the ranks}

There is also a finite closed form for the entire rank sequence. Let $P$ be
 the prime support of $k$, put $k_T=\prod_{p\in P\setminus T}p^{a_p}$, and set
\[
 C_\sigma(h)=2^h,\qquad C_\tau(h)=h+1,\qquad
 G_S(z)=\sum_{T\supseteq S}(-1)^{|T\setminus S|}
                  \frac{(k_T-1)z}{1-k_Tz}.
\]
\begin{corollary}[Rank generating function]\label{cor:rankseries}
For $f=\sigma$ or $\tau$, the formal generating function of the cumulative
kernel ranks is rational and equals
\[
 \sum_{e\ge0}(\dim_\Q V^f_{k,e})z^e
 =\frac{1-z^{C_f(|P|)}+
       \sum_{S\subseteq P}(1-z^{C_f(|S|)})G_S(z)}{(1-z)^2}.
\]
\end{corollary}
\begin{proof}
At birth level $t$, the number of admissible $u$ with $P_{t,u}\supseteq T$
is $(k_T-1)k_T^{t-1}$. Indeed $u$ must be a multiple of
$\prod_{p\in T}p^{ta_p}$; after removing that factor, the range has length
$k_T^t$ and multiples of $k_T$ are excluded. Boolean inclusion--exclusion
therefore makes $G_S$ the birth-level generating function for classes with
$P_{t,u}=S$ exactly. A class of capacity $C$ contributes
$\min(e-t+1,C)$ at level $e\ge t$, whose generating function is
$z^t(1-z^C)/(1-z)^2$. The zero class contributes the same expression with
$t=0$ and capacity $C_f(|P|)$. Sum the class contributions.
\end{proof}
This rationality concerns the dimension sequence, not the arithmetic
coefficient generating function. The dimensions themselves are unbounded;
the finite expression describes their growth without asserting regularity
of the underlying arithmetic sequence.

\section{An exact obstruction to finite rational dilation closure}\label{sec:dilation}

Fix integers $b\ge2$ and $r\ge1$, and let $a_n$ be a bounded integer sequence.
Define the germs at zero
\[
 A(z)=\sum_{n\ge1}a_nz^n,\qquad
 F_{a,r}(z)=\sum_{n\ge1}\frac{a_nz^n}{(b^n-1)^r},\qquad
 (\mathcal D_bF)(z)=F(bz).
\]
Here $b$ is fixed and $z$ is an auxiliary numerator variable. This is not the
Mahler substitution $z\mapsto z^b$, and it is not the generating function
obtained by replacing $b^{-1}$ with $z$ everywhere.

Coefficientwise cancellation gives
\begin{equation}\label{eq:dilation}
 (\mathcal D_b-1)^rF_{a,r}=A.
\end{equation}
For constant weights and $r=1$, the forcing term is $z/(1-z)$. Duverney's
account of the classical Lambert-series Pad\'e method displays precisely this
rational first-order forcing equation~\cite[Eq.~(10)]{duverney}.

\begin{theorem}[Finite rational closure criterion]\label{thm:dilation}
For bounded integer $a_n$, the following are equivalent:
\begin{enumerate}[label=(\roman*),itemsep=0pt]
\item $a_n$ is eventually periodic;
\item $F_{a,r}$ satisfies a nontrivial finite inhomogeneous rational linear
 dilation equation
 \[\sum_{j=0}^J R_j(z)F_{a,r}(b^jz)=R(z),\qquad R_j,R\in\C(z);\]
\item $F_{a,r}$ belongs to a finite-dimensional $\C(z)$-vector space of
 meromorphic germs stable under $\mathcal D_b$.
\end{enumerate}
If these conditions fail, the circle $|z|=b^r$ is a natural boundary of
$F_{a,r}$, and $1,F_{a,r}(z),F_{a,r}(bz),\ldots$ are linearly independent over
$\C(z)$.
\end{theorem}
\begin{proof}
A bounded integer sequence takes values in a finite alphabet. Its generating
function $A$ is rational if and only if it is eventually periodic. For the
nontrivial direction, a rational generating function gives an eventual
constant-coefficient recurrence. A state consisting of its finitely many
preceding coefficients takes only finitely many values, and the recurrence
determines the next state. A repeated state forces eventual periodicity.

Suppose $a_n$ is not eventually periodic. Infinitely many coefficients are
nonzero, so $A$ has radius one. The P\'olya--Carlson theorem gives a natural
boundary at the unit circle~\cite{bell-miles-ward}. Also $F_{a,r}$ has radius
$b^r$. Expanding \eqref{eq:dilation}, the term with largest dilation is
$F_{a,r}(b^rz)$, with coefficient one. Every other term is analytic on
$|z|<b$. Analytic continuation of $F_{a,r}$ through any point of $|z|=b^r$
would therefore continue $A$ through the corresponding unit-circle point,
a contradiction.

Now suppose a rational dilation relation exists, and let $J$ be the largest
index with $R_J\ne0$. Choose a generic point $z_0$ on
$|z|=b^{r-J}$ avoiding the finitely many poles of the rational coefficients
and zeros of $R_J$. All terms $F_{a,r}(b^jz)$ with $j<J$ are analytic near
$z_0$. Solving for $F_{a,r}(b^Jz)$ continues it through that point, contradicting
the natural boundary. A relation containing only the rational term is also
impossible unless that term is zero. This proves the final independence
assertion and excludes (ii).

If $a_n$ is eventually periodic, $A\in\Q(z)$ and \eqref{eq:dilation}
proves (ii). Moreover the span of
$1,F_{a,r},\mathcal D_bF_{a,r},\ldots,\mathcal D_b^{r-1}F_{a,r}$ is
finite-dimensional and stable under $\mathcal D_b$, since rational functions
remain rational on replacing $z$ by $bz$. This proves (iii). Conversely,
finite dimension forces a linear dependence among the dilates, giving (ii).
The identically zero and eventually zero sequences are included in this
periodic case.
\end{proof}

\begin{corollary}[M\"obius obstruction and a rational control]\label{cor:mobius}
For every $b\ge2$ and $r\ge1$, the M\"obius-weighted function $F_{\mu,r}$ has
natural boundary $|z|=b^r$ and admits no nontrivial finite rational linear
dilation equation. Nevertheless
\[
 F_{\mu,1}(1)=\frac1b,\qquad
 F_{\mu,2}(1)=\sum_{n\ge1}\frac{\ph(n)}{b^n}-\frac1b.
\]
\end{corollary}
\begin{proof}
If $\mu$ had an eventual period $T$, choose a prime $q\nmid T$. Every residue
class modulo $T$ contains arbitrarily large multiples of $q^2$, so the
periodic value on each class would have to be zero. This contradicts the
existence of arbitrarily large primes. Apply Theorem~\ref{thm:dilation}.

Absolute convergence allows geometric expansion and regrouping. For the
first power, the coefficient of $b^{-m}$ is $\sum_{d\mid m}\mu(d)$, which is
one at $m=1$ and zero otherwise. For the second power it is
\[
 \sum_{d\mid m}\mu(d)(m/d-1)
 =\ph(m)-\mathbf1_{m=1}.
\]
These give the two identities.
\end{proof}

A cross-problem comparison separates this analytic obstruction from arithmetic
rationality. Let $a_n$ be the indicator of the squares. It is not eventually
periodic, since its zero blocks have unbounded length and it has infinitely
many ones. The same natural-boundary and non-closure conclusion holds, but
\[
 F_{a,1}(1)=\sum_{j\ge1}\frac1{b^{j^2}-1}\notin\Q.
\]
This irrationality is a case of Duverney's Theorem~1~\cite{duverney} and of the
reciprocal-summable support theorem in the attached \#257 note. Thus the
M\"obius and square-support examples have the same analytic obstruction and
opposite arithmetic outcomes. For \#1049, the constant-weight Lambert family
lies on the finite-closure side instead.

The finite-alphabet arithmetic hypothesis is essential. Bounded rational
weights without a common denominator need not satisfy the dichotomy: take
$a_n=(1-b^{-n})^r$. These weights are nonperiodic, but
\[
 F_{a,r}(z)=\sum_{n\ge1}b^{-rn}z^n=\frac{z}{b^r-z}
\]
is rational. Bounded rational weights with a common denominator reduce to the
integer theorem by scaling.

The theorem excludes a finite rational linear closure mechanism in this
auxiliary variable. It does not exclude all Pad\'e approximants, nonlinear
functional equations, infinite auxiliary systems, or a value-specific
construction with effective denominator cancellation. The rational value at
$r=1$ is an explicit falsifier of any inference from the natural boundary or
dilation rank alone to irrationality. For $r=2,b=2$, the remaining arithmetic
quantity for rational approximants is still
\[
 0<q_N\left|F_{\mu,2}(1)-p_N/q_N\right|\longrightarrow0
\]
with $p_N/q_N$ in lowest terms. A determinant nonvanishing theorem alone does
not supply this denominator-scaled error estimate.

\section{Formalisation and falsification boundaries}

The three Lean files are uncompiled proof candidates for periodic transport,
Boolean inversion and coefficient-level cancellation. They do not formalise
the complete signature theorem, the divisor ranks or the analytic arguments.

The exact checker passes twenty $(f,k,e)$ cases over each of
$\mathbb F_{1000003}$ and $\mathbb F_{1000033}$, on both full channel matrices
and the proposed basis submatrices. All ranks match, including $230$ versus
$229$ at $(6,3)$. Four rational Boolean inverses, $450$ dilation coefficient
identities and $56$ coefficients of the rank generating function also pass.
These finite checks are reproducible; the universal conclusions depend on
the ordinary proofs, which remain subject to independent review.

\begin{thebibliography}{9}\small
\bibitem{martin} G.~Martin, \emph{Simultaneous inequalities among values of the
Euler phi-function}, arXiv:math/0603053, 2006. Theorem~1 and Corollary~4.
\bibitem{mattila} M.~Mattila and P.~Haukkanen, \emph{On the positive
definiteness and eigenvalues of meet and join matrices}, arXiv:1209.4287.
\bibitem{bell-miles-ward} J.~Bell, R.~Miles and T.~Ward, \emph{Towards a
P\'olya--Carlson dichotomy for algebraic dynamics}, arXiv:1307.2369.
The introduction states the classical integer-coefficient radius-one theorem.
\bibitem{duverney} D.~Duverney, \emph{Arithmetical functions and irrationality
of Lambert series}, author manuscript, Eq.~(10), p.~3.
\url{https://danielduverney.fr/documents/theorie-des-nombres/TokyoNT.pdf}.
\end{thebibliography}
\end{document}
